%%%%%%%%%%%%%%%%%%%%%%%%%%%%%%%%%%%%%%%%%%%%%%%%%%%%%%%%%%%%%%%%%%%%%%%%%%%%%%%%
% Introduction to Machine Learning (SS 2026) -- Programming Project
% Powerpredict: predicting daily power consumption from weather data.
%
% This is a SHARED report. Author 1 (k-NN, Method 1) content is filled in below.
% Author 2 should add the Ridge-Regression (Method 2) parts where marked
% "%% >>> RIDGE (Author 2)". Keep the whole report to 2 pages (excl. references).
%
% Items still to fill in are marked \TODO{...} (red). See the accompanying chat
% message / figure snippet for how to obtain each number.
%%%%%%%%%%%%%%%%%%%%%%%%%%%%%%%%%%%%%%%%%%%%%%%%%%%%%%%%%%%%%%%%%%%%%%%%%%%%%%%%

\documentclass[a4, 10 pt, conference]{ieeeconf}  % Comment this line out if you need a4paper

%\documentclass[a4paper, 10pt, conference]{ieeeconf}      % Use this line for a4 paper

\IEEEoverridecommandlockouts                              % This command is only needed if
                                                          % you want to use the \thanks command

\overrideIEEEmargins                                      % Needed to meet printer requirements.

\usepackage{multicol}
\usepackage{tcolorbox}
\usepackage{cuted,tcolorbox,lipsum}
\usepackage{xcolor}
\usepackage{graphicx}   % for \includegraphics (figures)
\usepackage{booktabs}   % nicer tables (\toprule, \midrule, \bottomrule)

% Helper for open items; delete all \TODO{...} before submitting.
\newcommand{\TODO}[1]{\textcolor{red}{[TODO: #1]}}

\title{\LARGE \bf
Introduction to Machine Learning (SS 2026)\\ Programming Project
\vspace{-3em}
}


\begin{document}


\maketitle
\vspace{-3em}
\thispagestyle{empty}
\pagestyle{empty}

\begin{strip}
\begin{tcolorbox}[
size=tight,
colback=white,
boxrule=0.2mm,
left=3mm,right=3mm, top=3mm, bottom=1mm
]
{\begin{multicols}{2}% replace 3 with 2 for 2 authors.

\textbf{Author 1}       \\
Last name: \TODO{last name}   \\
First name: \TODO{first name} \\
Matrikel Nr.: \TODO{Matrikel Nr.} \\

\columnbreak

\textbf{Author 2}       \\
Last name: Beneduce             \\
First name: Paul             \\
Matrikel Nr.: 12309055               \\

\end{multicols}}
\end{tcolorbox}
\end{strip}

%%%%%%%%%%%%%%%%%%%%%%%%%%%%%%%%%%%%%%%%%%%%%%%%%%%%%%%%%%%%%%%%%%%%%%%%%%%%%%%%

\section{Introduction}
\label{sec:intro}

We address a supervised \emph{regression} task: predicting the daily total
\texttt{power\_consumption} of a region from weather observations. Accurately
anticipating demand from weather is directly relevant for energy suppliers, who
must match generation to load.

The dataset (\texttt{powerpredict}) contains $39{,}997$ instances and $65$
input features. Each feature describes the weather in one of five fictional cities
(Bedrock, Gotham City, New New York, Paperopoli, Springfield). Per city there
are $11$ numeric measurements (temperature and its daily low/high, pressure,
humidity, wind speed and direction, short-/long-term rain, snow, cloud cover)
and $2$ textual weather descriptors (\texttt{weather\_main},
\texttt{weather\_description}), giving $55$ numeric and $10$ categorical columns
in total. The target is a single continuous value. The numeric features contain
no missing values, so no imputation was required, and they live on very
different scales (e.g. pressure vs.\ rainfall), which motivates the scaling step
described below. The target is \TODO{describe target distribution from
eda.ipynb: e.g. right-skewed, range and mean $\approx$ value}. We evaluate all
models with the \emph{Mean Absolute Error} (MAE), the official competition
metric (lower is better).

\section{Implementation / ML Process}
\label{sec:methods}

\subsection*{Data preprocessing}
Both methods share one preprocessing stage, implemented as a scikit-learn
\texttt{ColumnTransformer} inside a \texttt{Pipeline}. Numeric columns are
standardised with \texttt{StandardScaler} (zero mean, unit variance). This is
essential for a distance-based method such as k-NN: without it, large-range
features (e.g.\ pressure) would dominate the neighbour search and drown out
small-range ones (e.g.\ rainfall). The textual weather descriptors are turned
into indicator columns with \texttt{OneHotEncoder}
(\texttt{handle\_unknown="ignore"}, so unseen weather categories in the hidden
test set do not cause errors). We treated the use of these categorical columns
as a design choice and compared both variants (Section~\ref{sec:results}).
Wrapping preprocessing and model in a single pipeline guarantees that the scaler
and encoder are fitted only on the training folds, preventing data leakage
during cross-validation, and lets the evaluation notebook feed in the raw table
directly.

\subsection*{Method 1: k-Nearest Neighbors regression}
Method~1 is a k-Nearest-Neighbors regressor. It is an instance-based,
non-parametric method: it stores the (scaled) training data and predicts a new
point by averaging the targets of its $k$ closest neighbours under the Euclidean
(Minkowski, $p=2$) distance. We use distance weighting, so closer neighbours
contribute more strongly to the prediction than distant ones.

k-NN is a natural fit here: the weather-to-power relationship is likely
non-linear and locally smooth (similar weather days yield similar consumption),
which is exactly the structure k-NN exploits without assuming a global
functional form. It also serves as a strong, simple contrast to the linear
Method~2. As a non-parametric estimator it is closely related to the
non-parametric density estimation covered in the VO~\cite{vo-nonparam}; k-NN
regression replaces the density estimate with a local average of the target.

\subsection*{Hyperparameters}
We selected hyperparameters with an exhaustive grid search using
cross-validation, scored by (negative) MAE, so the choice reflects performance
on held-out folds rather than the training data. The search grid was the number
of neighbours $k \in \{3, 5, 11, 25, 51\}$ and the weighting scheme
(\emph{uniform} or \emph{distance}). The best configuration was
$k = 51$ with distance weighting. The final submitted model uses the numeric
features only (categorical encoding disabled): dropping the high-cardinality
one-hot columns kept memory within the JupyterHub limit and did not meaningfully
hurt accuracy (Section~\ref{sec:results}).

%% >>> RIDGE (Author 2): add a subsection here describing Ridge regression
%% (regularised linear regression), why it suits the task, and how you chose
%% the regularisation strength alpha via cross-validation.

\section{Results}
\label{sec:results}

Table~\ref{tab:cv} reports the cross-validated MAE of the k-NN model across the
hyperparameter grid, and Fig.~\ref{fig:knn} visualises how the validation MAE
changes with $k$. Larger $k$ with distance weighting gave the lowest error, i.e.
$k=51$ was selected. Table~\ref{tab:final} summarises the final performance
against the dummy baseline (a regressor that always predicts the training-set
mean).

\begin{table}[h]
\centering
\caption{Cross-validated MAE of the k-NN model for each hyperparameter
configuration (lower is better).}
\label{tab:cv}
\begin{tabular}{ccc}
\toprule
$k$ (neighbours) & weights & CV MAE \\
\midrule
3   & uniform  & \TODO{} \\
3   & distance & \TODO{} \\
5   & uniform  & \TODO{} \\
5   & distance & \TODO{} \\
11  & uniform  & \TODO{} \\
11  & distance & \TODO{} \\
25  & uniform  & \TODO{} \\
25  & distance & \TODO{} \\
51  & uniform  & \TODO{} \\
\textbf{51} & \textbf{distance} & \textbf{3509.56} \\
\bottomrule
\end{tabular}
\end{table}

\begin{figure}[h]
\centering
% Generate figures/knn_val_mae_vs_k.png with the snippet provided, then uncomment:
% \includegraphics[width=0.9\linewidth]{figures/knn_val_mae_vs_k.png}
\TODO{insert figures/knn\_val\_mae\_vs\_k.png}
\caption{Validation MAE of the k-NN regressor as a function of the number of
neighbours $k$, for uniform and distance weighting.}
\label{fig:knn}
\end{figure}

\begin{table}[h]
\centering
\caption{Final MAE on the training, validation and (hidden) test sets. Lower is
better.}
\label{tab:final}
\begin{tabular}{lccc}
\toprule
Model & Train MAE & Valid.\ MAE & Test MAE \\
\midrule
Dummy (mean)      & \TODO{} & \TODO{} & \TODO{} \\
k-NN (Method 1)   & 15.69   & 3509.56 & 2438.15 \\
%% >>> RIDGE (Author 2): add Ridge row here.
Ridge (Method 2)  & \TODO{} & \TODO{} & \TODO{} \\
\bottomrule
\end{tabular}
\end{table}

\section{Discussion}
\label{sec:discuss}

The k-NN model clearly beats the dummy baseline on unseen data
(test MAE $2438$ vs.\ dummy \TODO{value}), confirming that the weather features
carry real predictive signal. The very low \emph{training} MAE ($15.69$) is
expected and not a sign of quality: with distance weighting, a query that
coincides with a stored training point receives essentially that point's own
target (weight $\to \infty$ at distance $0$), so the model reproduces the
training data almost exactly. The gap to the validation/test MAE ($\approx 2400$--$3500$)
is the true generalisation error and the number that matters.

Interestingly, the test MAE ($2438$) is lower than the cross-validation estimate
($3509$). This is plausible because each CV fold trains on only two thirds of
the data, whereas the final model is refit on the full training set, and because
the hidden test set may be slightly easier; in any case it indicates the model
is not overfitting to the public data.

Two things we tried but dropped: (i) one-hot encoding the weather-text columns
inflated the feature space and, together with the memory footprint of k-NN,
exhausted the JupyterHub RAM (the training process was killed); the gain in
accuracy was negligible \TODO{cite the with/without-encoding numbers from
training.ipynb, cell~4}, so we kept the leaner numeric-only model. (ii) Running
the grid search in parallel (\texttt{n\_jobs=-1}) replicated the full dataset
across worker processes and triggered the same out-of-memory failure, so the
search is run sequentially.

The best $k$ ($51$) sat at the upper edge of our grid, which suggests even larger
neighbourhoods might reduce error further; extending the grid (e.g.\ $k \in
\{75, 101\}$) is the obvious next step. Beyond that, dimensionality reduction
such as PCA~\cite{vo-pca} could counter the curse of dimensionality that hurts
distance-based methods, and feature selection or an alternative distance metric
could sharpen the neighbour search.

\section{Conclusion}
\label{sec:con}

Our final k-NN regressor achieves a test-set MAE of $\mathbf{2438.15}$,
substantially better than the dummy baseline of \TODO{value}. We did not tune
against the test set; hyperparameters were fixed purely on cross-validated
training performance.

The main takeaway is that a simple, non-parametric method with careful
preprocessing goes a long way on this problem: standardising the weather
features and letting k-NN average over locally similar days already yields a
large improvement over the mean predictor. The exercise also highlighted the
practical trade-offs of instance-based learning: k-NN's memory cost forced us to
drop the categorical features and run the search sequentially, and its tendency
to memorise the training data makes the training error misleading. A regularised
linear model (Method~2) offers a complementary, low-memory perspective, and
comparing the two clarifies when local averaging beats a global linear fit for
weather-driven demand prediction.

%% >>> RIDGE (Author 2): extend the conclusion with Method 2's test performance
%% and a one-line comparison of the two methods.

%%%%%%%%%%%%%%%%%%%%%%%%%%%%%%%%%%%%%%%%%%%%%%%%%%%%%%%%%%%%%%%%%%%%%%%%%%%%%%%%

\begin{thebibliography}{9}

\bibitem{vo-nonparam}
Introduction to Machine Learning (SS 2026), Lecture~10:
\emph{Nonparametric Density Estimation}. University of Innsbruck.

\bibitem{vo-pca}
Introduction to Machine Learning (SS 2026), Lecture~07:
\emph{Dimensionality Reduction by PCA/SVD}. University of Innsbruck.

\end{thebibliography}

\end{document}
